\section{Metodología}

Este artículo se desarrolló bajo un enfoque cualitativo de investigación documental, orientado a la revisión y análisis de literatura científica relacionada con paradigmas, metodologías y arquitecturas de software. De acuerdo con Metodología Scrum: Desarrollo detallado de la fase de aprobación de un proyecto informático mediante el uso de metodologías ágiles \cite{scrum}, los enfoques cualitativos permiten interpretar fenómenos complejos y aportar lineamientos prácticos basados en evidencia académica.

La realización de este proyecto se centró en comprender y aplicar principios clave como los de la Programación Orientada a Objetos \cite{poo}, la arquitectura Modelo–Vista–Controlador \cite{mvc}, la gestión de bases de datos \cite{databases} y el uso de metodologías ágiles como Scrum \cite{scrum}. Estos lineamientos técnicos y metodológicos habían sido definidos desde las etapas iniciales del proceso formativo como parte del plan de desarrollo de un sistema de gestión de eventos. La revisión documental permitió sustentar dichas decisiones y evidenciar sus beneficios en términos de organización, colaboración y robustez en el desarrollo de software.

Sus fases fueron:

\subsection{Fase 1- Recolección y construcción del thesaurus}

En esta primera etapa se buscó consolidar un sustento teórico para el proyecto. Para ello, se recurrió a Google Scholar como principal fuente de información, aplicando búsquedas con términos clave vinculados a metodologías ágiles, patrones de diseño, Programación Orientada a Objetos (POO) y el modelo MVC. Los resultados más pertinentes se registraron en un documento tipo thesaurus, en el que se organizaron de manera uniforme datos como el título del artículo, la fuente de publicación, la fecha y el concepto central. El objetivo de esta fase fue recolectar y estructurar evidencias académicas que demostraran la pertinencia de emplear las metodologías y tecnologías seleccionadas dentro del desarrollo del proyecto formativo.

\subsection{Fase 2- Análisis cualitativo de la evidencia (síntesis del thesaurus)}

La finalidad de esta etapa consistió en convertir la información recopilada en sustento teórico y práctico que respaldara la incorporación de enfoques como el Modelo–Vista–Controlador (MVC), la Programación Orientada a Objetos (POO) y la metodología ágil Scrum en el desarrollo del sistema. Para ello, se llevaron a cabo procesos de lectura analítica, elaboración de resúmenes y contrastación crítica de los artículos seleccionados, lo que permitió identificar elementos clave como la separación de responsabilidades en la arquitectura, la facilidad de mantenimiento del software y el valor del trabajo colaborativo en equipos ágiles.

\subsection{Fase 3- Diseño y desarrollo del sistema}

El propósito central de esta fase fue llevar a la práctica los fundamentos teóricos y metodológicos previamente analizados, convirtiéndolos en un sistema de gestión de eventos plenamente funcional. Para ello, se tomó como base la arquitectura Modelo–Vista–Controlador (MVC), que facilitó la organización en capas; la Programación Orientada a Objetos (POO), que aportó principios de reutilización y escalabilidad; y la metodología ágil Scrum, que orientó el proceso de construcción mediante una gestión iterativa y colaborativa. En conjunto, estas guías permitieron estructurar y coordinar el proyecto de manera ordenada, asegurando tanto la calidad técnica del software como la efectividad en la organización del trabajo del equipo.